% English abstract; one page, structured, Times New Roman 12 pt
\ChapterStart
\begin{latin}
\begin{center}
  {\bfseries\fontsize{14pt}{17pt}\selectfont \ThesisTitleEN\par}
  \vspace{1em}
  {\bfseries\fontsize{14pt}{17pt}\selectfont Abstract}
\end{center}
\vspace{0.8em}
\noindent\textbf{Background and Aim:}
Hypertension, as a major chronic disease and an important risk factor for cardiovascular diseases, requires continuous self-care behaviors for effective management. Workshop-based education is beneficial because it provides opportunities for interaction and hands-on practice; however, its time-limited nature may reduce the long-term retention of educational content. The present study aimed to compare the effectiveness of workshop-based self-care education with workshop-based education integrated with microlearning on self-care performance and blood pressure among patients with hypertension.

\vspace{0.45em}
\noindent\textbf{Methods:}
This randomized controlled trial was conducted among 110 patients with hypertension attending the cardiology clinic of Shahid Madani Hospital in Khorramabad, Iran, during the first half of 1405 in the Solar Hijri calendar. Participants were assigned to two groups—workshop-based education and workshop-based education integrated with microlearning—using block randomization. Both groups received four 60-minute face-to-face educational sessions at weekly intervals. In addition to these sessions, the integrated intervention group received eight 3-minute educational microvideos. Self-care performance was assessed using the Persian version of the Hypertension Self-Care Performance Questionnaire, and systolic and diastolic blood pressure indices were assessed using measurements recorded in a blood pressure log. Assessments were performed at two time points: before the intervention and one month after the intervention. Data were analyzed using SPSS software and paired-samples t-tests, independent-samples t-tests, and analysis of covariance (ANCOVA).

\vspace{0.45em}
\noindent\textbf{Results:}
A total of 47 participants were included in each group. The mean total self-care score increased from 59.70 to 86.89 in the workshop-based education integrated with microlearning group and from 62.23 to 81.55 in the workshop-based education group. After adjustment for baseline scores, the adjusted between-group difference was 5.61 points in favor of the workshop-based education integrated with microlearning group (95\% confidence interval [CI]: 2.03–9.19; p = 0.002). The adjusted between-group differences in the information, control, and awareness domains were also statistically significant in favor of the workshop-based education integrated with microlearning group; however, no significant differences were observed between the groups in the attitude and behavior domains. Systolic and diastolic blood pressure decreased following the intervention in both groups; however, these reductions were not statistically significant.

\vspace{0.45em}
\noindent\textbf{Conclusion:}
Adding video-based microlearning to workshop-based educational sessions resulted in a significant improvement in self-care performance, particularly in the information, control, and awareness domains, among patients with hypertension. However, compared with conventional workshop-based education, this combined approach did not demonstrate a significant advantage in reducing systolic or diastolic blood pressure. Therefore, microlearning may be used as an effective complementary educational strategy alongside face-to-face instruction to facilitate the repetition and reinforcement of educational messages and to enhance patients’ self-care knowledge and skills.

\vspace{0.8em}
\noindent\textbf{Keywords:} \KeywordsEN
\end{latin}
\clearpage
